
\section{Graph Theory}

 A \emph{Graph} \(G = \inner{V}{E} \) consists of a finite set of vertices \(V\) and a relation 
 \(E \subseteq V \times V\) over the set. Depending on the traits of \(E\) it can be categorized into 

 \begin{itemize}
    
    \item \emph{Undirected}: \(E\) is symmetric.

    \item \emph{Directed}: \(E\) is not symmetric.

    \item \emph{Complete}: \(E = V \times V\) for directed and \(|E| \approx |V| (|V| - 1) / 2\) for undirected graphs.
    
    \item \emph{Loop-free}: \(\forall(v,v) \in V \times V : (v,v) \notin E \)

    \item \emph{Weighted}: \(G = \langle V, E, \phi \rangle\), where \(\phi: E \to R\) and \(R\) is an arbitrary set.
    It mostly refers to adding a cost to each edge.
    
    \item \emph{Dense}: \(|E| \approx |V|^2\) for directed and \(|E| \approx |V| (|V| - 1) / 2\) for undirected graphs.

    \item \emph{Sparse}: \(|E| << |V|^2\) for directed and \(|E| << |V| (|V| - 1) / 2\) for undirected graphs.

 \end{itemize}

 In computer science the density of a graph is a main criteria when modeling a graph. 

 \subsection{Connected Graphs}

 A graph \(G = (V,E)\) undirected, then it is called \emph{connected} if for all pairs \((v_i, v_j) \subseteq V \times V\) there exists a 
 path from \(v_i\) to \(v_j\). We define \(E_S\) the symmetric relation. 

 \[
    (v_1, v_2) \in E \iff (v_1, v_2) \in E_S \land (v_2, v_1) \in E_S
 \]

 If \(G\) is directed then it is connected if and only if its undirected version is connected.
 
 \subsection{Paths}

 Given a graph \(G\), we call a \emph{path} a set of vertices if 

 \[
    \forall i \in \{1, \dots, k - 1\} : (w_i, w_{i + 1}) \in E
 \]

 The most common types of paths are 

\begin{itemize}
    
    \item \emph{Circle}: \(w_1 = w_k\)

    \item \emph{Eulerian Circuit}: Each edge will be visited one time. \(\forall(s,t) \in E : \exists_1 i: s = w_i \land t = w_{i + 1}\)

    \item \emph{Hamiltonian Circuit}: Each vertex will be visited one time \(\forall v \in V : \exists_1 i : v = w_i\) 
    
    \item \emph{Loop-free}: \(\forall(v,v) \in V \times V : (v,v) \notin E \)

 \end{itemize}

 \subsubsection{Path Cost}

 We define the \emph{path cost} for a weighted graph as 

 \[
    \phi (p) = \sum_{i = 1}^{k - 1} \phi(w_i, w_{i+1})
 \]

 \subsection{The Adjacency Matrix}

 For a graph \(G\) with vertices \(V = \{v_1, \dots, v_n\}\), we define the \emph{the adjacency matrix} \(A = (a_{ij})\) as a 
 \(n \times n\) matrix with 

 \[
    a_{ij} = \begin{cases}
        1 & if (v_i, v_j) \in E \\ 
        0 & else
    \end{cases}
 \]

\subsection{The Relax Operation}

In graph theory when referring to the action of updating the cost of visiting a vertex in a weigthed graph we call this operation 
\emph{relaxation}. 

\[
    \operatorname{relaz}(u,v, \text{weight}) = 
    \begin{cases}
        d[u] + \text{weight}(u,v) &\text{if} d[u] + \text{weight}(u,v) < d[v] \\ 
        d[u] & \text{else}
    \end{cases}
\]



